\documentclass[12pt,a4paper]{article}

% ---------------------------------------------------------
% PACKAGES
% ---------------------------------------------------------
\usepackage[margin=1in]{geometry}
\usepackage{amsmath, amssymb}
\usepackage{mathtools}
\usepackage{graphicx}
\usepackage{setspace}
\usepackage{hyperref}
\hypersetup{
    hidelinks
}
\usepackage{bm}
\usepackage{titlesec}
\usepackage{graphicx}
\usepackage{xcolor}
\usepackage{titling}
\usepackage[font=small,labelfont=bf]{caption}
\captionsetup{justification=centering}
\usepackage{titlesec}
\usepackage{listings}
\usepackage{xcolor}

\titleformat{\section}
  {\normalfont\Large\bfseries} 
  {\thesection.}
  {0.75em}
  {}

\titleformat{\subsection}
  {\normalfont\large\bfseries}
  {\thesubsection.}
  {0.75em}
  {}

\titleformat{\subsubsection}
  {\normalfont\normalsize\bfseries}
  {\thesubsubsection.}
  {0.75em}
  {}



\titlelabel{\thetitle.\quad}

\setstretch{1.2}

% ---------------------------------------------------------
\lstdefinestyle{mystyle}{
  language=Matlab,
  basicstyle=\ttfamily\small,
  keywordstyle=\color{blue},
  commentstyle=\color{gray},
  numbers=left,
  numberstyle=\tiny,
  stepnumber=1,
  frame=single,
  breaklines=true,
  captionpos=b
}




\begin{document}
% ---------------- Cover Page ----------------
\begin{titlepage}
\centering

\vspace*{1.8cm}

{\LARGE \textbf{Lebanese American University}\par}
\vspace{1cm}

% ---- LAU Logo (replace lau_logo.png with your uploaded file name) ----
\includegraphics[width=0.55\textwidth]{Photos/lau.png}\par
\vspace{1.4cm}

{\large \textbf{ELE 443 -- Control Systems Lab}\par}
\vspace{0.2cm}
{\large Section 31\par}
\vspace{0.2cm}
{\large Mr. Ziad Saoud\par}

\vspace{0.8cm}

{\Large \textbf{Furuta Pendulum Project Report}\par}
\vspace{0.3cm}
{\large Fall 2025\par}

\vfill

% ---- Bottom right block (names/IDs/date) ----
\begin{flushright}
\textbf{Ali Awada} \hspace{0.5em} 202300797\\
\textbf{Mahdi T. Zein Al Dine} \hspace{0.5em} 202300041\\[6pt]
15/12/2025
\end{flushright}

\end{titlepage}

% ---------------- Cover Page ----------------



% ---------------- Tables ----------------
\tableofcontents
\newpage

\listoffigures
\newpage

% ---------------------------------------------------------
% SECTIONS
% ---------------------------------------------------------

\section{Introduction}
The Furuta pendulum, also known as the rotary inverted pendulum, is a classical benchmark system in control engineering due to its nonlinear dynamics, underactuated structure, and inherent instability around the upright equilibrium. It provides a practical platform for studying modeling techniques, nonlinear dynamics, linearization, state-space analysis, and advanced control strategies within a laboratory environment.

This project aims to develop a complete mathematical and control-oriented description of the Furuta pendulum system, starting from first-principle modeling and progressing toward controller design and performance evaluation. The system is modeled using Lagrangian mechanics to capture the coupled dynamics between the rotary arm and the pendulum link. From the derived equations of motion, linearized models are obtained around both the downward stable equilibrium and the upright unstable equilibrium.

Based on the linearized models, state-space representations are constructed and used to analyze system behavior and controllability. Multiple control approaches are then explored, including open-loop simulations, PID control, and full state feedback control. Each control strategy is simulated and evaluated to assess its effectiveness in stabilizing the pendulum and regulating the system response.

The overall objective of this report is to bridge theoretical control concepts with practical implementation, highlighting the challenges associated with nonlinear systems, equilibrium-dependent behavior, and real-world constraints. Through systematic modeling, simulation, and control design, the Furuta pendulum serves as a comprehensive educational tool for advanced control systems analysis.

\section{Modeling, Physical Parameters, and Assumptions}
In this section, the physical structure of the Furuta pendulum is described along with the parameters and assumptions used to develop its mathematical model. The goal of the modeling process is to capture the dominant dynamics of the system while maintaining a tractable formulation suitable for analysis, linearization, and control design. All parameters are chosen based on the laboratory setup and Datasheet specifications, and simplifying assumptions are introduced where appropriate to reduce model complexity without significantly affecting accuracy.

\subsection*{Modeling Assumptions}

To derive the equations of motion, the following assumptions are made:

\begin{itemize}
    \item All mechanical components are rigid bodies.
    \item The rotary arm and pendulum are connected by ideal revolute joints.
    \item Friction at the joints is modeled as viscous damping.
    \item Air resistance and structural flexibility are neglected.
    \item The pendulum mass is concentrated at its center of mass.
    \item Gravitational acceleration is constant.
\end{itemize}

These assumptions are standard in Furuta pendulum modeling and provide a reasonable balance between physical realism and mathematical simplicity.

\subsection*{Physical Parameters and Coordinates}

The system consists of a horizontal rotary arm driven by a DC motor and a pendulum link attached at the end of the arm. The generalized coordinates are defined as:
\begin{itemize}
    \item $\theta$: rotation angle of the rotary arm about the vertical axis
    \item $\alpha$: angular displacement of the pendulum measured from the vertical direction
\end{itemize}

A schematic of the Furuta pendulum with annotated parameters and coordinate definitions is shown in Fig.~\ref{fig:furuta_pendulum}.

\begin{figure}[h!]
    \centering
    \includegraphics[width=0.4\textwidth]{Photos/AnnotatedPendelom.png}
    \caption{Annotated Furuta pendulum }
    \label{fig:furuta_pendulum}
\end{figure}

% =========================================================
\section{Lagrangian Formulation}

\subsection{Assumptions, Parameters, and Coordinates}

\textbf{Given physical parameters:}
\begin{itemize}
  \item Rotary arm: mass $m_r$, length $r$, inertia about vertical axis through pivot $j_r$, viscous damping $b_r$.
  \item Pendulum link: mass $m_p$, total length $L_p$, center-of-mass distance $h_g=\dfrac{L_p}{2}$, inertia about its center of mass $j_p$, viscous damping $b_p$.
  \item Gravity: $g$.
\end{itemize}

\textbf{Generalized coordinates:}
\[
q_1=\theta(t), \qquad q_2=\alpha(t).
\]
\begin{itemize}
  \item $\theta$ is the rotary arm angle about the vertical axis ($+Z$), positive for counter-clockwise (CCW) rotation.
  \item $\alpha$ is the pendulum angle measured from the upright vertical direction:
  \[
  \alpha=0 \;\Rightarrow\; \text{pendulum perfectly upright}, \qquad \alpha\neq 0 \;\Rightarrow\; \text{rotation about the arm axis}.
  \]
\end{itemize}

\textbf{Actuation and dissipation:}
\begin{itemize}
  \item Motor $V_m$ applies torque $\tau_m$ about $\theta$.
  \item Viscous friction is modeled by :
  \[
  \mathcal{F}=\frac{1}{2}b_r\dot\theta^2+\frac{1}{2}b_p\dot\alpha^2.
  \]
\end{itemize}

\subsection{Reference Frames and Kinematics}

Define an inertial frame $\{\bm e_X,\bm e_Y,\bm e_Z\}$ with $\bm e_Z$ upward.

Define arm-attached unit vectors in the horizontal plane:
\[
\bm e_r =
\begin{bmatrix}\cos\theta\\ \sin\theta\\ 0\end{bmatrix},
\qquad
\bm e_t =
\begin{bmatrix}-\sin\theta\\ \cos\theta\\ 0\end{bmatrix},
\qquad
\bm e_Z=
\begin{bmatrix}0\\ 0\\ 1\end{bmatrix}.
\]
These satisfy the time-derivatives:
\[
\dot{\bm e}_r=\dot\theta\,\bm e_t,
\qquad
\dot{\bm e}_t=-\dot\theta\,\bm e_r,
\qquad
\dot{\bm e}_Z=\bm 0.
\]

\subsubsection{Position of the Pendulum Center of Mass}

The rotary arm end-point (pendulum pivot) is
\[
\bm r_A = r\,\bm e_r.
\]

The pendulum swings in a vertical plane containing $\bm e_t$ and $\bm e_Z$, and rotates about the arm axis $\bm e_r$.
Choose the unit vector along the pendulum (from pivot toward the COM) as
\[
\bm u(\alpha)= -\sin\alpha\,\bm e_t+\cos\alpha\,\bm e_Z,
\]
so that:
\[
\alpha=0 \Rightarrow \bm u=\bm e_Z \quad (\text{upright}).
\]

Hence the COM position is
\[
\bm r_G
=
\bm r_A+h_g\,\bm u
=
r\,\bm e_r+h_g\big(-\sin\alpha\,\bm e_t+\cos\alpha\,\bm e_Z\big).
\]

\subsubsection{Velocity of the Pendulum Center of Mass}

Differentiate $\bm r_G$:
\[
\dot{\bm r}_G
=
r\,\dot{\bm e}_r
+h_g\left(-\cos\alpha\,\dot\alpha\,\bm e_t-\sin\alpha\,\dot{\bm e}_t-\sin\alpha\,\dot\alpha\,\bm e_Z\right).
\]

Substitute $\dot{\bm e}_r=\dot\theta\,\bm e_t$ and $\dot{\bm e}_t=-\dot\theta\,\bm e_r$:
\begin{align*}
\dot{\bm r}_G
&=
r\,\dot\theta\,\bm e_t
+h_g\left(-\cos\alpha\,\dot\alpha\,\bm e_t-\sin\alpha(-\dot\theta\,\bm e_r)-\sin\alpha\,\dot\alpha\,\bm e_Z\right)\\[4pt]
&=
\big(r\dot\theta-h_g\cos\alpha\,\dot\alpha\big)\bm e_t
+\big(h_g\sin\alpha\,\dot\theta\big)\bm e_r
+\big(-h_g\sin\alpha\,\dot\alpha\big)\bm e_Z.
\end{align*}

Therefore the squared speed is
\begin{align*}
v_G^2
&=\dot{\bm r}_G\cdot\dot{\bm r}_G\\
&=\big(r\dot\theta-h_g\cos\alpha\,\dot\alpha\big)^2
+\big(h_g\sin\alpha\,\dot\theta\big)^2
+\big(h_g\sin\alpha\,\dot\alpha\big)^2\\[4pt]
&=
\left(r^2+h_g^2\sin^2\alpha\right)\dot\theta^2
-2r h_g\cos\alpha\,\dot\theta\dot\alpha
+h_g^2\dot\alpha^2.
\end{align*}

\subsubsection{Angular Velocity of the Pendulum Link (About its COM)}

The link experiences:
\begin{itemize}
  \item yaw rate $\dot\theta$ about $\bm e_Z$,
  \item swing rate $\dot\alpha$ about $\bm e_r$.
\end{itemize}

Thus the angular velocity is
\[
\bm\omega_p=\dot\theta\,\bm e_Z+\dot\alpha\,\bm e_r.
\]

With the common rigid-body approximation used in Furuta modeling,
the rotational kinetic energy about the COM is taken as
\[
T_{\text{rot},p}=\frac{1}{2}j_p\left(\dot\alpha^2+\dot\theta^2\sin^2\alpha\right).
\]

\subsection{Energies}

\subsubsection{Kinetic Energy}

\textbf{Rotary arm kinetic energy} (about the vertical axis):
\[
T_r=\frac{1}{2}j_r\dot\theta^2.
\]

\textbf{Pendulum translational kinetic energy}:
\[
T_{\text{trans},p}=\frac{1}{2}m_p v_G^2
=\frac{1}{2}m_p\left[\left(r^2+h_g^2\sin^2\alpha\right)\dot\theta^2
-2r h_g\cos\alpha\,\dot\theta\dot\alpha
+h_g^2\dot\alpha^2\right].
\]

\textbf{Pendulum rotational kinetic energy} (about COM):
\[
T_{\text{rot},p}=\frac{1}{2}j_p\left(\dot\alpha^2+\dot\theta^2\sin^2\alpha\right).
\]

\textbf{Total kinetic energy}:
\begin{align*}
T
&=T_r+T_{\text{trans},p}+T_{\text{rot},p}\\[4pt]
&=\frac{1}{2}\left(j_r+m_p r^2+m_p h_g^2\sin^2\alpha+j_p\sin^2\alpha\right)\dot\theta^2
-\left(m_p r h_g\cos\alpha\right)\dot\theta\dot\alpha
+\frac{1}{2}\left(j_p+m_p h_g^2\right)\dot\alpha^2.
\end{align*}

For compactness in later steps, define:
\[
J_t \coloneqq j_r+m_p r^2,
\qquad
J_\alpha \coloneqq j_p+m_p h_g^2.
\]
Then:
\[
T=\frac{1}{2}\left(J_t+J_\alpha\sin^2\alpha\right)\dot\theta^2
-\left(m_p r h_g\cos\alpha\right)\dot\theta\dot\alpha
+\frac{1}{2}J_\alpha\dot\alpha^2.
\]

\subsubsection{Potential Energy}

The COM height is the $Z$-component of $\bm r_G$:
\[
z_G=\bm r_G\cdot\bm e_Z=h_g\cos\alpha.
\]
Thus the gravitational potential energy is
\[
V=m_p g z_G=m_p g h_g\cos\alpha.
\]

\subsubsection{Lagrangian}

\[
\mathcal{L}=T-V.
\]
So explicitly:
\begin{equation}
\boxed{
\mathcal{L}
=
\frac{1}{2}\left(J_t+J_\alpha\sin^2\alpha\right)\dot\theta^2
-\left(m_p r h_g\cos\alpha\right)\dot\theta\dot\alpha
+\frac{1}{2}J_\alpha\dot\alpha^2
-\;m_p g h_g\cos\alpha
}
\label{eq:lagrangain}
\end{equation}

% =========================================================
% =========================================================
\section{Equations of Motion}

Starting from the Lagrangian
\begin{equation}
\mathcal{L}
=
\frac{1}{2}\left(J_t+J_\alpha\sin^2\alpha\right)\dot\theta^2
-\left(m_p r h_g\cos\alpha\right)\dot\theta\dot\alpha
+\frac{1}{2}J_\alpha\dot\alpha^2
-\;m_p g h_g\cos\alpha,
\label{eq:lagrangian}
\end{equation}
the equations of motion are derived using the Euler--Lagrange equations with viscous damping:
\begin{equation}
\frac{d}{dt}\left(\frac{\partial\mathcal{L}}{\partial\dot q_i}\right)
-\frac{\partial\mathcal{L}}{\partial q_i}
+\frac{\partial\mathcal{F}}{\partial\dot q_i}
=Q_i,
\qquad
(q_1,q_2)=(\theta,\alpha),
\label{eq:eulerlagrange}
\end{equation}
where the friction dissipation function is
\begin{equation}
\mathcal{F}=\frac{1}{2}b_r\dot\theta^2+\frac{1}{2}b_p\dot\alpha^2,
\label{eq:dissipation}
\end{equation}
and the generalized forces are
\[
(Q_\theta,Q_\alpha)=(\tau_m,0).
\]

% ---------------------------------------------------------
\subsection{Equation of Motion in $\theta$}

\paragraph{Partial derivative with respect to $\theta$.}

The Lagrangian \eqref{eq:lagrangian} does not depend explicitly on $\theta$, hence
\begin{equation}
\frac{\partial\mathcal{L}}{\partial\theta}=0.
\label{eq:dLdtheta}
\end{equation}

\paragraph{Partial derivative with respect to $\dot\theta$.}

From \eqref{eq:lagrangian},
\begin{align}
\frac{\partial\mathcal{L}}{\partial\dot\theta}
&=
\left(J_t+J_\alpha\sin^2\alpha\right)\dot\theta
-\left(m_p r h_g\cos\alpha\right)\dot\alpha.
\label{eq:dLdthetadot}
\end{align}

\paragraph{Time derivative of $\partial\mathcal{L}/\partial\dot\theta$.}

Taking the time derivative of \eqref{eq:dLdthetadot}:
\begin{align}
\frac{d}{dt}\left(\frac{\partial\mathcal{L}}{\partial\dot\theta}\right)
&=
\left(J_t+J_\alpha\sin^2\alpha\right)\ddot\theta
+J_\alpha\sin(2\alpha)\dot\alpha\dot\theta
\nonumber\\
&\quad
-\left(m_p r h_g\right)\frac{d}{dt}\left(\cos\alpha\,\dot\alpha\right).
\label{eq:d_dtdLdthetadot}
\end{align}

The remaining derivative is
\begin{align}
\frac{d}{dt}\left(\cos\alpha\,\dot\alpha\right)
&=
-\sin\alpha\,\dot\alpha^2+\cos\alpha\,\ddot\alpha.
\label{eq:cosalphadot}
\end{align}

Substituting \eqref{eq:cosalphadot} into \eqref{eq:d_dtdLdthetadot} yields
\begin{align}
\frac{d}{dt}\left(\frac{\partial\mathcal{L}}{\partial\dot\theta}\right)
&=
\left(J_t+J_\alpha\sin^2\alpha\right)\ddot\theta
+J_\alpha\sin 2\alpha\,\dot\theta\dot\alpha
\nonumber\\
&\quad
+m_p r h_g\sin\alpha\,\dot\alpha^2
-\;m_p r h_g\cos\alpha\,\ddot\alpha.
\label{eq:d_dtdLdthetadot_final}
\end{align}

\paragraph{Damping term.}

From \eqref{eq:dissipation},
\begin{equation}
\frac{\partial\mathcal{F}}{\partial\dot\theta}=b_r\dot\theta.
\label{eq:dFdthetadot}
\end{equation}

\paragraph{Final equation in $\theta$.}

Substituting \eqref{eq:dLdtheta}, \eqref{eq:d_dtdLdthetadot_final}, and \eqref{eq:dFdthetadot} into \eqref{eq:eulerlagrange} gives:
\begin{equation}
\boxed{
\left(J_t+J_\alpha\sin^2\alpha\right)\ddot\theta
-\;m_p r h_g\cos\alpha\,\ddot\alpha
+J_\alpha\sin 2\alpha\,\dot\theta\dot\alpha
+m_p r h_g\sin\alpha\,\dot\alpha^2
+b_r\dot\theta
=
\tau_m.
}
\label{eq:lagrangiantheta}
\end{equation}

% ---------------------------------------------------------
\subsection{Equation of Motion in $\alpha$}

\paragraph{Partial derivative with respect to $\alpha$.}

Differentiating \eqref{eq:lagrangian} with respect to $\alpha$:
\begin{align}
\frac{\partial\mathcal{L}}{\partial\alpha}
&=
J_\alpha\sin\alpha\cos\alpha\,\dot\theta^2
+m_p r h_g\sin\alpha\,\dot\theta\dot\alpha
+m_p g h_g\sin\alpha.
\label{eq:dLdalpha}
\end{align}

\paragraph{Partial derivative with respect to $\dot\alpha$.}

From \eqref{eq:lagrangian},
\begin{equation}
\frac{\partial\mathcal{L}}{\partial\dot\alpha}
=
-\left(m_p r h_g\cos\alpha\right)\dot\theta
+J_\alpha\dot\alpha.
\label{eq:dLdalphadot}
\end{equation}

\paragraph{Time derivative of $\partial\mathcal{L}/\partial\dot\alpha$.}

Taking the time derivative of \eqref{eq:dLdalphadot}:
\begin{align}
\frac{d}{dt}\left(\frac{\partial\mathcal{L}}{\partial\dot\alpha}\right)
&=
-\left(m_p r h_g\right)\frac{d}{dt}\left(\cos\alpha\,\dot\theta\right)
+J_\alpha\ddot\alpha.
\label{eq:d_dtdLdalphadot}
\end{align}

The remaining derivative is
\begin{align}
\frac{d}{dt}\left(\cos\alpha\,\dot\theta\right)
&=
-\sin\alpha\,\dot\alpha\dot\theta+\cos\alpha\,\ddot\theta.
\label{eq:cosalphathetadot}
\end{align}

Substituting \eqref{eq:cosalphathetadot} into \eqref{eq:d_dtdLdalphadot} yields
\begin{align}
\frac{d}{dt}\left(\frac{\partial\mathcal{L}}{\partial\dot\alpha}\right)
&=
J_\alpha\ddot\alpha
-\;m_p r h_g\cos\alpha\,\ddot\theta
+m_p r h_g\sin\alpha\,\dot\theta\dot\alpha.
\label{eq:d_dtdLdalphadot_final}
\end{align}

\paragraph{Damping term.}

From \eqref{eq:dissipation},
\begin{equation}
\frac{\partial\mathcal{F}}{\partial\dot\alpha}=b_p\dot\alpha.
\label{eq:dFdalphadot}
\end{equation}

\paragraph{Final equation in $\alpha$.}

Substituting \eqref{eq:dLdalpha}, \eqref{eq:d_dtdLdalphadot_final}, and \eqref{eq:dFdalphadot} into \eqref{eq:eulerlagrange} gives:
\begin{equation}
\boxed{
J_\alpha\ddot\alpha
-\;m_p r h_g\cos\alpha\,\ddot\theta
-\;\frac{J_\alpha}{2}\sin 2\alpha\,\dot\theta^2
-\;m_p g h_g\sin\alpha
+b_p\dot\alpha
=
0.
}
\label{eq:lagrangianalpha}
\end{equation}

% ---------------------------------------------------------
\subsection{Final Nonlinear Equations of Motion}

Recalling
\[
J_t=j_r+m_p r^2,
\qquad
J_\alpha=j_p+m_p h_g^2,
\]
the nonlinear equations of motion of the Furuta pendulum are given by
\eqref{eq:lagrangiantheta} and \eqref{eq:lagrangianalpha}.
% =========================================================

% =========================================================
\section{Linearization Around the Upright Equilibrium ($\alpha=0$)}

This section linearizes the nonlinear equations of motion around the upright equilibrium
\[
\alpha_0 = 0, \qquad \theta=\theta_0, \qquad \dot\alpha_0=0,\qquad \dot\theta_0=0,
\]
(with $\theta_0$ arbitrary; without loss of generality we may take $\theta_0=0$ since $\theta$ does not appear explicitly in the dynamics).

% ---------------------------------------------------------
\subsection{Nonlinear EOM (starting point)}

Recall the nonlinear equations of motion:
\begin{align}
\left(J_t+J_\alpha\sin^2\alpha\right)\ddot\theta
-\;m_p r h_g\cos\alpha\,\ddot\alpha
+J_\alpha\sin 2\alpha\,\dot\theta\dot\alpha
+m_p r h_g\sin\alpha\,\dot\alpha^2
+b_r\dot\theta
&=
\tau_m,
\label{eq:eom_theta_nl}\\[6pt]
J_\alpha\ddot\alpha
-\;m_p r h_g\cos\alpha\,\ddot\theta
-\;\frac{J_\alpha}{2}\sin 2\alpha\,\dot\theta^2
-\;m_p g h_g\sin\alpha
+b_p\dot\alpha
&=
0.
\label{eq:eom_alpha_nl}
\end{align}
where
\[
J_t=j_r+m_p r^2,\qquad J_\alpha=j_p+m_p h_g^2.
\]

Define for compactness:
\[
B \;\coloneqq\; m_p r h_g.
\]

% ---------------------------------------------------------
\subsection{Small-angle approximations about $\alpha=0$}

For small deviations around $\alpha=0$:
\[
\sin\alpha \approx \alpha,\qquad
\cos\alpha \approx 1,\qquad
\sin^2\alpha \approx 0,\qquad
\sin 2\alpha \approx 2\alpha.
\]

Additionally, since we linearize about the rest condition $(\dot\theta_0,\dot\alpha_0)=(0,0)$,
any product of small quantities such as
\[
\alpha\,\dot\theta\,\dot\alpha,\qquad
\alpha\,\dot\alpha^2,\qquad
\alpha\,\dot\theta^2
\]
is higher-order and is neglected in the linear model.

% ---------------------------------------------------------
\subsection{Linearization of the $\theta$-equation}

Start from \eqref{eq:eom_theta_nl}. Apply the approximations:
\[
J_t+J_\alpha\sin^2\alpha \approx J_t,\qquad
\cos\alpha \approx 1,
\]
and drop the nonlinear velocity-product terms:
\[
J_\alpha\sin 2\alpha\,\dot\theta\dot\alpha \approx 2J_\alpha\,\alpha\,\dot\theta\dot\alpha \;\approx\; 0,
\qquad
B\sin\alpha\,\dot\alpha^2 \approx B\,\alpha\,\dot\alpha^2 \;\approx\; 0.
\]

Thus, the linearized $\theta$ equation becomes:
\begin{equation}
\boxed{
J_t\,\ddot\theta \;-\; B\,\ddot\alpha \;+\; b_r\,\dot\theta \;=\; \tau_m.
}
\label{eq:eom_theta_lin0}
\end{equation}

% ---------------------------------------------------------
\subsection{Linearization of the $\alpha$-equation}

Start from \eqref{eq:eom_alpha_nl}. Apply:
\[
\cos\alpha \approx 1,\qquad
\sin\alpha \approx \alpha,
\]
and drop the nonlinear term:
\[
-\frac{J_\alpha}{2}\sin 2\alpha\,\dot\theta^2 \approx -J_\alpha\,\alpha\,\dot\theta^2 \;\approx\; 0.
\]

Then:
\begin{equation}
\boxed{
J_\alpha\,\ddot\alpha \;-\; B\,\ddot\theta \;-\; m_p g h_g\,\alpha \;+\; b_p\,\dot\alpha \;=\; 0.
}
\label{eq:eom_alpha_lin0}
\end{equation}

% ---------------------------------------------------------
\subsection{Linearized model in matrix form}

Group \eqref{eq:eom_theta_lin0}--\eqref{eq:eom_alpha_lin0} as:
\begin{equation}
\boxed{
\begin{bmatrix}
J_t & -B\\
-B & J_\alpha
\end{bmatrix}
\begin{bmatrix}
\ddot\theta\\
\ddot\alpha
\end{bmatrix}
=
\begin{bmatrix}
\tau_m - b_r\dot\theta\\
m_p g h_g\,\alpha - b_p\dot\alpha
\end{bmatrix}.
}
\label{eq:lin0_matrix_form}
\end{equation}

Define the determinant:
\[
\Delta \coloneqq J_tJ_\alpha - B^2.
\]

% ---------------------------------------------------------
\subsection{Explicit expressions for $\ddot\theta$ and $\ddot\alpha$}

Inverting the $2\times2$ inertia matrix in \eqref{eq:lin0_matrix_form} gives:
\begin{align}
\boxed{
\ddot\theta
=
\frac{1}{\Delta}
\left[
J_\alpha\left(\tau_m-b_r\dot\theta\right)
+
B\left(m_p g h_g\,\alpha-b_p\dot\alpha\right)
\right],
}
\label{eq:theta_ddot_lin0}\\[6pt]
\boxed{
\ddot\alpha
=
\frac{1}{\Delta}
\left[
B\left(\tau_m-b_r\dot\theta\right)
+
J_t\left(m_p g h_g\,\alpha-b_p\dot\alpha\right)
\right].
}
\label{eq:alpha_ddot_lin0}
\end{align}

\noindent\textit{Remark (stability):}
Since the gravity term appears as $+m_p g h_g\,\alpha$ on the right-hand side of
\eqref{eq:alpha_ddot_lin0}, the upright equilibrium $\alpha=0$ is open-loop unstable.

% (State-space is placed in the next report section.)
% =========================================================


% =========================================================
\section{State-Space Representation}

This section converts the linearized model around the upright equilibrium $\alpha=0$ into a first-order state-space form. We start from the linearized second-order equations obtained in the previous section and then define the state vector, derive each state derivative explicitly, and finally assemble the matrices $A,B,C,D$.

% ---------------------------------------------------------
\subsection{Starting point: linearized equations about $\alpha=0$}

From \eqref{eq:eom_theta_lin0}--\eqref{eq:eom_alpha_lin0}, the linearized dynamics are:
\begin{align}
J_t\,\ddot\theta \;-\; B\,\ddot\alpha \;+\; b_r\,\dot\theta &= \tau_m,
\label{eq:lin0_theta_2nd}\\
J_\alpha\,\ddot\alpha \;-\; B\,\ddot\theta \;-\; m_p g h_g\,\alpha \;+\; b_p\,\dot\alpha &= 0,
\label{eq:lin0_alpha_2nd}
\end{align}
with the compact constants
\[
B \coloneqq m_p r h_g,
\qquad
\Delta \coloneqq J_t J_\alpha - B^2.
\]

Rearrange \eqref{eq:lin0_theta_2nd}--\eqref{eq:lin0_alpha_2nd} into matrix form:
\begin{equation}
\begin{bmatrix}
J_t & -B\\
-B & J_\alpha
\end{bmatrix}
\begin{bmatrix}
\ddot\theta\\
\ddot\alpha
\end{bmatrix}
=
\begin{bmatrix}
\tau_m - b_r\dot\theta\\
m_p g h_g\,\alpha - b_p\dot\alpha
\end{bmatrix}.
\label{eq:lin0_Mqdd}
\end{equation}

% ---------------------------------------------------------
\subsection{Solving explicitly for $\ddot\theta$ and $\ddot\alpha$}

The inverse of a $2\times2$ matrix
$
\begin{bmatrix} a & b\\ c & d \end{bmatrix}
$
is
$
\dfrac{1}{ad-bc}
\begin{bmatrix} d & -b\\ -c & a \end{bmatrix}
$.
Here, in \eqref{eq:lin0_Mqdd},
\[
a=J_t,\quad b=-B,\quad c=-B,\quad d=J_\alpha,
\qquad
ad-bc = J_tJ_\alpha - (-B)(-B)=J_tJ_\alpha-B^2=\Delta.
\]

Therefore,
\[
\begin{bmatrix}
\ddot\theta\\
\ddot\alpha
\end{bmatrix}
=
\frac{1}{\Delta}
\begin{bmatrix}
J_\alpha & B\\
B & J_t
\end{bmatrix}
\begin{bmatrix}
\tau_m - b_r\dot\theta\\
m_p g h_g\,\alpha - b_p\dot\alpha
\end{bmatrix}.
\]

Carrying out the multiplication gives the explicit accelerations:
\begin{align}
\ddot\theta
&=
\frac{1}{\Delta}
\Big[
J_\alpha(\tau_m-b_r\dot\theta)
+
B(m_p g h_g\,\alpha-b_p\dot\alpha)
\Big],
\label{eq:theta_ddot_lin0_again}\\[6pt]
\ddot\alpha
&=
\frac{1}{\Delta}
\Big[
B(\tau_m-b_r\dot\theta)
+
J_t(m_p g h_g\,\alpha-b_p\dot\alpha)
\Big].
\label{eq:alpha_ddot_lin0_again}
\end{align}

% ---------------------------------------------------------
\subsection{State vector definition}

Choose the standard 4-state vector:
\[
\bm x \coloneqq
\begin{bmatrix}
x_1\\x_2\\x_3\\x_4
\end{bmatrix}
=
\begin{bmatrix}
\theta\\ \alpha\\ \dot\theta\\ \dot\alpha
\end{bmatrix},
\qquad
u \coloneqq \tau_m.
\]

Then, by definition:
\[
\dot x_1 = x_3,
\qquad
\dot x_2 = x_4.
\]

The remaining two derivatives are:
\[
\dot x_3 = \ddot\theta,
\qquad
\dot x_4 = \ddot\alpha,
\]
which we now rewrite using \eqref{eq:theta_ddot_lin0_again}--\eqref{eq:alpha_ddot_lin0_again}.

% ---------------------------------------------------------
\subsection{Deriving each state derivative explicitly}

\subsubsection{$\dot x_1$ and $\dot x_2$}

\[
\boxed{\dot x_1 = x_3,}
\qquad
\boxed{\dot x_2 = x_4.}
\]

\subsubsection{$\dot x_3$}

From \eqref{eq:theta_ddot_lin0_again} with $u=\tau_m$, $\dot\theta=x_3$, $\alpha=x_2$, $\dot\alpha=x_4$:
\begin{align*}
\dot x_3
&=
\frac{1}{\Delta}\Big[J_\alpha(u-b_r x_3)+B(m_p g h_g\,x_2-b_p x_4)\Big]\\[4pt]
&=
\underbrace{\left(\frac{B\,m_p g h_g}{\Delta}\right)}_{\text{coeff of }x_2}x_2
+\underbrace{\left(-\frac{J_\alpha b_r}{\Delta}\right)}_{\text{coeff of }x_3}x_3
+\underbrace{\left(-\frac{B b_p}{\Delta}\right)}_{\text{coeff of }x_4}x_4
+\underbrace{\left(\frac{J_\alpha}{\Delta}\right)}_{\text{coeff of }u}u.
\end{align*}

So:
\[
\boxed{
\dot x_3
=
\frac{B m_p g h_g}{\Delta}x_2
-\frac{J_\alpha b_r}{\Delta}x_3
-\frac{B b_p}{\Delta}x_4
+\frac{J_\alpha}{\Delta}u.
}
\]

\subsubsection{$\dot x_4$}

From \eqref{eq:alpha_ddot_lin0_again} with $u=\tau_m$, $\dot\theta=x_3$, $\alpha=x_2$, $\dot\alpha=x_4$:
\begin{align*}
\dot x_4
&=
\frac{1}{\Delta}\Big[B(u-b_r x_3)+J_t(m_p g h_g\,x_2-b_p x_4)\Big]\\[4pt]
&=
\underbrace{\left(\frac{J_t\,m_p g h_g}{\Delta}\right)}_{\text{coeff of }x_2}x_2
+\underbrace{\left(-\frac{B b_r}{\Delta}\right)}_{\text{coeff of }x_3}x_3
+\underbrace{\left(-\frac{J_t b_p}{\Delta}\right)}_{\text{coeff of }x_4}x_4
+\underbrace{\left(\frac{B}{\Delta}\right)}_{\text{coeff of }u}u.
\end{align*}

So:
\[
\boxed{
\dot x_4
=
\frac{J_t m_p g h_g}{\Delta}x_2
-\frac{B b_r}{\Delta}x_3
-\frac{J_t b_p}{\Delta}x_4
+\frac{B}{\Delta}u.
}
\]

% ---------------------------------------------------------
\subsection{Assembling the state-space matrices}

Collecting
$\dot{\bm x} = A\bm x + B u$ with $\bm x=[\theta,\alpha,\dot\theta,\dot\alpha]^T$:

\[
\boxed{
\dot{\bm x}
=
\underbrace{
\begin{bmatrix}
0 & 0 & 1 & 0\\
0 & 0 & 0 & 1\\
0 & \dfrac{B m_p g h_g}{\Delta} & -\dfrac{J_\alpha b_r}{\Delta} & -\dfrac{B b_p}{\Delta}\\
0 & \dfrac{J_t m_p g h_g}{\Delta} & -\dfrac{B b_r}{\Delta} & -\dfrac{J_t b_p}{\Delta}
\end{bmatrix}
}_{A}
\bm x
+
\underbrace{
\begin{bmatrix}
0\\
0\\
\dfrac{J_\alpha}{\Delta}\\
\dfrac{B}{\Delta}
\end{bmatrix}
}_{B}
u.
}
\]

% ---------------------------------------------------------
\subsection{Output equation $y=Cx+Du$}

For the output to be the angles of the arm and the pendulum, we may define:
\[
\bm y \coloneqq
\begin{bmatrix}
\theta\\
\alpha
\end{bmatrix},
\]
then
\[
\boxed{
\bm y
=
\underbrace{
\begin{bmatrix}
1 & 0 & 0 & 0\\
0 & 1 & 0 & 0
\end{bmatrix}
}_{C}
\bm x
+
\underbrace{
\begin{bmatrix}
0\\
0
\end{bmatrix}
}_{D}
u.
}
\]

% ---------------------------------------------------------
\subsection{Final state-space model}

With
$
\bm x=[\theta,\alpha,\dot\theta,\dot\alpha]^T
$
and
$
u=\tau_m
$,
the linearized system around $\alpha=0$ is:
\[
\boxed{
\dot{\bm x}=A\bm x+Bu,\qquad \bm y=C\bm x+Du,
}
\]
where $A,B,C,D$ are given explicitly above, and
\[
B=m_p r h_g,\qquad \Delta=J_tJ_\alpha-B^2.
\]
% ---------------------------------------------------------
\subsection{Conversion from Motor Torque $\tau_m$ to Motor Voltage $V_m$}

In practice, the Furuta pendulum is actuated by a DC motor whose input is the armature voltage $V_m$, not the motor torque $\tau_m$. Therefore, the state-space model derived previously must be reformulated to use $V_m$ as the control input.

\subsubsection{DC motor model}

For a DC motor operating in the linear region, the electrical and mechanical relations are:
\begin{align}
\tau_m &= K_t i, \label{eq:tau_current}\\
V_m &= R_m i + K_m \dot\theta, \label{eq:voltage_eq}
\end{align}
where:
\begin{itemize}
  \item $i$ is the armature current,
  \item $K_t$ is the motor torque constant,
  \item $K_m$ is the back-EMF constant,
  \item $R_m$ is the armature resistance.
\end{itemize}

Solving \eqref{eq:voltage_eq} for the current:
\[
i = \frac{V_m - K_m\dot\theta}{R_m}.
\]

Substitute into \eqref{eq:tau_current}:
\begin{equation}
\boxed{
\tau_m
=
\frac{K_t}{R_m}V_m
-
\frac{K_t K_m}{R_m}\dot\theta.
}
\label{eq:tau_voltage}
\end{equation}

Thus, the motor torque consists of a voltage-dependent actuation term and an additional velocity-dependent damping term.

% ---------------------------------------------------------
\subsubsection{Modified state equations with voltage input}

Recall that in the torque-input state-space model,
the torque $\tau_m$ appears only in the $\dot x_3$ and $\dot x_4$ equations.
Substitute \eqref{eq:tau_voltage} into those equations.

\paragraph{Equation for $\dot x_3$:}

From the torque-input model,
\[
\dot x_3
=
\frac{B m_p g h_g}{\Delta}x_2
-\frac{J_\alpha b_r}{\Delta}x_3
-\frac{B b_p}{\Delta}x_4
+\frac{J_\alpha}{\Delta}\tau_m.
\]

Substituting $\tau_m$ from \eqref{eq:tau_voltage}:
\[
\dot x_3
=
\frac{B m_p g h_g}{\Delta}x_2
-\frac{J_\alpha}{\Delta}\!\left(b_r+\frac{K_t K_m}{R_m}\right)x_3
-\frac{B b_p}{\Delta}x_4
+\frac{J_\alpha K_t}{\Delta R_m}V_m.
\]

\paragraph{Equation for $\dot x_4$:}

Similarly, from the torque-input model,
\[
\dot x_4
=
\frac{J_t m_p g h_g}{\Delta}x_2
-\frac{B b_r}{\Delta}x_3
-\frac{J_t b_p}{\Delta}x_4
+\frac{B}{\Delta}\tau_m.
\]

Substituting $\tau_m$:
\[
\dot x_4
=
\frac{J_t m_p g h_g}{\Delta}x_2
-\frac{B}{\Delta}\!\left(b_r+\frac{K_t K_m}{R_m}\right)x_3
-\frac{J_t b_p}{\Delta}x_4
+\frac{B K_t}{\Delta R_m}V_m.
\]

% ---------------------------------------------------------
\subsubsection{State-space matrices with voltage input}

Define the new input as
\[
u \coloneqq V_m.
\]

Then the voltage-driven state-space model becomes:
\[
\boxed{
\dot{\bm x}
=
A_v \bm x
+
B_v u,
}
\]
where
\[
A_v =
\begin{bmatrix}
0 & 0 & 1 & 0\\
0 & 0 & 0 & 1\\
0 & \dfrac{B m_p g h_g}{\Delta}
& -\dfrac{J_\alpha}{\Delta}\!\left(b_r+\dfrac{K_t K_m}{R_m}\right)
& -\dfrac{B b_p}{\Delta}\\
0 & \dfrac{J_t m_p g h_g}{\Delta}
& -\dfrac{B}{\Delta}\!\left(b_r+\dfrac{K_t K_m}{R_m}\right)
& -\dfrac{J_t b_p}{\Delta}
\end{bmatrix},
\]

\[
B_v =
\begin{bmatrix}
0\\
0\\
\dfrac{J_\alpha K_t}{\Delta R_m}\\
\dfrac{B K_t}{\Delta R_m}
\end{bmatrix}.
\]

The output equation remains unchanged:
\[
\boxed{
\bm y = C\bm x + D u.
}
\]


% =========================================================
\section{Linearization Around the Downward Equilibrium ($\alpha=\pi$)}

In this section, the nonlinear equations of motion are linearized about the \emph{downward} equilibrium
\[
\alpha_0=\pi,\qquad \dot\alpha_0=0,\qquad \dot\theta_0=0,
\]
which corresponds to the pendulum hanging down. This equilibrium is (locally) stable in open-loop.

% ---------------------------------------------------------
\subsection{Nonlinear EOM (starting point)}

We start from the nonlinear EOM:
\begin{align}
\left(J_t+J_\alpha\sin^2\alpha\right)\ddot\theta
-\;m_p r h_g\cos\alpha\,\ddot\alpha
+J_\alpha\sin 2\alpha\,\dot\theta\dot\alpha
+m_p r h_g\sin\alpha\,\dot\alpha^2
+b_r\dot\theta
&=
\tau_m,
\label{eq:eom_theta_nl_pi}\\[6pt]
J_\alpha\ddot\alpha
-\;m_p r h_g\cos\alpha\,\ddot\theta
-\;\frac{J_\alpha}{2}\sin 2\alpha\,\dot\theta^2
-\;m_p g h_g\sin\alpha
+b_p\dot\alpha
&=
0,
\label{eq:eom_alpha_nl_pi}
\end{align}
with
\[
J_t=j_r+m_p r^2,\qquad J_\alpha=j_p+m_p h_g^2,\qquad B\coloneqq m_p r h_g.
\]

% ---------------------------------------------------------
\subsection{Deviation variable about $\alpha=\pi$}

Define the small deviation from the downward equilibrium:
\[
\beta(t)\;\coloneqq\;\alpha(t)-\pi,
\]
so that
\[
\alpha=\pi+\beta,\qquad \dot\alpha=\dot\beta,\qquad \ddot\alpha=\ddot\beta,
\qquad |\beta|\ll 1.
\]

Use the Taylor approximations around $\beta=0$:
\[
\sin(\pi+\beta)\approx-\beta,
\qquad
\cos(\pi+\beta)\approx-1,
\qquad
\sin^2(\pi+\beta)=\sin^2\beta\approx 0,
\qquad
\sin(2(\pi+\beta))=\sin(2\beta)\approx 2\beta.
\]
Also, products of small quantities such as $\beta\,\dot\theta\,\dot\beta$, $\beta\,\dot\theta^2$, and $\beta\,\dot\beta^2$ are neglected in the linear model.

% ---------------------------------------------------------
\subsection{Linearization of the $\theta$-equation}

Start from \eqref{eq:eom_theta_nl_pi}. Apply:
\[
J_t+J_\alpha\sin^2\alpha \approx J_t,\qquad \cos\alpha\approx-1,\qquad \sin\alpha\approx-\beta,
\]
and drop the nonlinear velocity-product terms:
\[
J_\alpha\sin2\alpha\,\dot\theta\dot\alpha \approx 2J_\alpha\beta\,\dot\theta\dot\beta \approx 0,
\qquad
B\sin\alpha\,\dot\alpha^2 \approx -B\beta\,\dot\beta^2 \approx 0.
\]

Therefore:
\begin{equation}
\boxed{
J_t\,\ddot\theta \;+\; B\,\ddot\beta \;+\; b_r\,\dot\theta \;=\; \tau_m.
}
\label{eq:linpi_theta_2nd_down}
\end{equation}

% ---------------------------------------------------------
\subsection{Linearization of the $\alpha$-equation}

Start from \eqref{eq:eom_alpha_nl_pi}. Apply:
\[
\cos\alpha\approx-1,\qquad \sin\alpha\approx-\beta,\qquad \sin2\alpha\approx2\beta,
\]
and drop the nonlinear term
\[
-\frac{J_\alpha}{2}\sin2\alpha\,\dot\theta^2 \approx -J_\alpha\beta\,\dot\theta^2 \approx 0.
\]

Then:
\begin{align*}
0
&=
J_\alpha\ddot\beta
-\;B\cos(\pi+\beta)\,\ddot\theta
-\;m_p g h_g\sin(\pi+\beta)
+b_p\dot\beta\\
&\approx
J_\alpha\ddot\beta
-\;B(-1)\ddot\theta
-\;m_p g h_g(-\beta)
+b_p\dot\beta,
\end{align*}
so the linearized $\beta$-equation is:
\begin{equation}
\boxed{
J_\alpha\,\ddot\beta \;+\; B\,\ddot\theta \;+\; m_p g h_g\,\beta \;+\; b_p\,\dot\beta \;=\; 0.
}
\label{eq:linpi_beta_2nd_down}
\end{equation}

% ---------------------------------------------------------
\subsection{Linear model in matrix form}

Group \eqref{eq:linpi_theta_2nd_down}--\eqref{eq:linpi_beta_2nd_down} as:
\begin{equation}
\boxed{
\begin{bmatrix}
J_t & B\\
B & J_\alpha
\end{bmatrix}
\begin{bmatrix}
\ddot\theta\\
\ddot\beta
\end{bmatrix}
=
\begin{bmatrix}
\tau_m - b_r\dot\theta\\
-\,m_p g h_g\,\beta - b_p\dot\beta
\end{bmatrix}.
}
\label{eq:linpi_matrix_down}
\end{equation}

Define the determinant:
\[
\Delta \coloneqq J_tJ_\alpha - B^2.
\]

% ---------------------------------------------------------
\subsection{Explicit expressions for $\ddot\theta$ and $\ddot\beta$}

Invert the $2\times2$ inertia matrix in \eqref{eq:linpi_matrix_down}:
\[
\begin{bmatrix}
\ddot\theta\\
\ddot\beta
\end{bmatrix}
=
\frac{1}{\Delta}
\begin{bmatrix}
J_\alpha & -B\\
-B & J_t
\end{bmatrix}
\begin{bmatrix}
\tau_m - b_r\dot\theta\\
-\,m_p g h_g\,\beta - b_p\dot\beta
\end{bmatrix}.
\]

Hence:
\begin{align}
\boxed{
\ddot\theta
=
\frac{1}{\Delta}\Big[
J_\alpha(\tau_m-b_r\dot\theta)
+
B\big(m_p g h_g\,\beta + b_p\dot\beta\big)
\Big],
}
\label{eq:theta_ddot_linpi_down}\\[6pt]
\boxed{
\ddot\beta
=
\frac{1}{\Delta}\Big[
-\,B(\tau_m-b_r\dot\theta)
-\;J_t\big(m_p g h_g\,\beta + b_p\dot\beta\big)
\Big].
}
\label{eq:beta_ddot_linpi_down}
\end{align}

\noindent\textit{Remark (stability):}
The gravity term appears with a restoring sign in \eqref{eq:beta_ddot_linpi_down}, so the downward equilibrium is locally stable in open-loop.

% =========================================================
\section{Simulation Around the Downward Equilibrium ($\alpha=\pi$)}

This section describes the simulation procedure and expected behavior when the system is initialized near the downward equilibrium. The simulation can be performed using either the torque-input model ($u=\tau_m$) or the voltage-input model ($u=V_m$).

% ---------------------------------------------------------
\subsection{State-space form about $\alpha=\pi$ (using $\beta=\alpha-\pi$)}

Define the state vector:
\[
\bm x_d \coloneqq
\begin{bmatrix}
x_1\\x_2\\x_3\\x_4
\end{bmatrix}
=
\begin{bmatrix}
\theta\\ \beta\\ \dot\theta\\ \dot\beta
\end{bmatrix},
\qquad
u\coloneqq\tau_m.
\]
Then:
\[
\dot x_1=x_3,\qquad \dot x_2=x_4,
\qquad
\dot x_3=\ddot\theta,\qquad \dot x_4=\ddot\beta.
\]

Substitute \eqref{eq:theta_ddot_linpi_down}--\eqref{eq:beta_ddot_linpi_down} with
$\dot\theta=x_3$, $\beta=x_2$, $\dot\beta=x_4$:
\begin{align*}
\dot x_3 &=
\frac{B m_p g h_g}{\Delta}x_2
-\frac{J_\alpha b_r}{\Delta}x_3
+\frac{B b_p}{\Delta}x_4
+\frac{J_\alpha}{\Delta}u,\\[4pt]
\dot x_4 &=
-\frac{J_t m_p g h_g}{\Delta}x_2
+\frac{B b_r}{\Delta}x_3
-\frac{J_t b_p}{\Delta}x_4
-\frac{B}{\Delta}u.
\end{align*}

Therefore,
\[
\boxed{\dot{\bm x}_d = A_d \bm x_d + B_d u}
\]
with
\[
A_d=
\begin{bmatrix}
0 & 0 & 1 & 0\\
0 & 0 & 0 & 1\\
0 & \dfrac{B m_p g h_g}{\Delta} & -\dfrac{J_\alpha b_r}{\Delta} & \dfrac{B b_p}{\Delta}\\
0 & -\dfrac{J_t m_p g h_g}{\Delta} & \dfrac{B b_r}{\Delta} & -\dfrac{J_t b_p}{\Delta}
\end{bmatrix},
\qquad
B_d=
\begin{bmatrix}
0\\
0\\
\dfrac{J_\alpha}{\Delta}\\
-\dfrac{B}{\Delta}
\end{bmatrix}.
\]

If the outputs are the angles:
\[
\bm y \coloneqq
\begin{bmatrix}\theta\\ \beta\end{bmatrix},
\qquad
C_d=
\begin{bmatrix}
1&0&0&0\\
0&1&0&0
\end{bmatrix},
\qquad
D_d=
\begin{bmatrix}0\\0\end{bmatrix}.
\]

% ---------------------------------------------------------
\subsection{Voltage-driven model about $\alpha=\pi$ (optional)}

Using the DC motor relation
\[
\tau_m=\frac{K_t}{R_m}V_m-\frac{K_tK_m}{R_m}\dot\theta,
\]
define the input $u\coloneqq V_m$. This substitution is equivalent to replacing
\[
b_r \;\mapsto\; b_r+\frac{K_tK_m}{R_m}
\]
in the $\dot\theta$ damping terms and scaling the input matrix by $\dfrac{K_t}{R_m}$:
\[
B_{d,v}=
\begin{bmatrix}
0\\
0\\
\dfrac{J_\alpha K_t}{\Delta R_m}\\
-\dfrac{B K_t}{\Delta R_m}
\end{bmatrix}.
\]
The corresponding $A_{d,v}$ is identical to $A_d$ except that every occurrence of $b_r$ is replaced by
$b_r+\dfrac{K_tK_m}{R_m}$.

% ---------------------------------------------------------
\subsection{Simulation setup}

The simulation is performed near the downward equilibrium using small initial deviations:
\[
\theta(0)=0,\quad \beta(0)=0,\quad \dot\theta(0)=0,\quad \dot\beta(0)=0,
\]

\subsection{MATLAB State-Space Implementation and Simulation Setup}

To validate the analytical model, the derived state-space representation is redefined and implemented directly in \textsc{MATLAB}. This allows systematic simulation using Simulink tools and ensures consistency with the numerical results obtained in later control design stages.

The physical parameters are first defined symbolically or numerically, after which the state-space matrices $(A,B,C,D)$ are constructed explicitly. The resulting linear model is then simulated using Simulink.

\begin{center}
\begin{lstlisting}[style=mystyle,
caption={MATLAB implementation of the state-space model around $\alpha=\pi$ (downward equilibrium)}]

clear; clc;

%% ===== Parameters (edit as needed) =====
mr  = 0.095;      % rotary arm mass [kg]
Lr  = 0.085;      % rotary arm length [m]
mp  = 0.024;      % pendulum mass [kg]
Lp  = 0.129;      % pendulum length [m]

Jm  = 1.4e-6;     % motor rotor inertia [kg*m^2]
Jh  = 0.6e-6;     % hub inertia [kg*m^2]

Rm  = 7.5;        % motor terminal resistance [Ohm]
Kt  = 0.0422;     % torque constant [N*m/A]
Km  = 0.0422;     % back-EMF constant [V/(rad/s)]

g   = 9.81;       % gravity [m/s^2]
br  = 0.001;      % arm viscous damping [N*m*s/rad]
bp  = 0.00005;    % pendulum viscous damping [N*m*s/rad]

%% ===== Derived geometry/inertias =====
r  = Lr;
hg = Lp/2;

jr = (1/3)*mr*Lr^2 ;    % effective rotary inertia about pivot`
jp = (1/3)*mp*Lp^2;              % pendulum inertia about its COM (as stated)

Jt     = jr ;
Jalpha = jp ;

B = mp*r*hg;   % coupling term

%% ===== Linearized equations in matrix form =====
% M * [th_dd; al_dd] = [ tau - br*th_d ;
%                       -bp*al_d - mp*g*hg*al ]
M = [ Jt,     B;
    B,  Jalpha ];

Minv = inv(M);

%% ===== State vector =====
% x = [ theta; alpha_tilde; theta_dot; alpha_dot ]
% u = tau_m  (torque input)
A = zeros(4,4);
B_tau = zeros(4,1);

% xdot1 = theta_dot
A(1,3) = 1;

% xdot2 = alpha_dot
A(2,4) = 1;


A(3,2) = Minv(1,2)*(-mp*g*hg);   % from -mp*g*hg*alpha
A(3,3) = Minv(1,1)*(-br);        % from -br*theta_dot
A(3,4) = Minv(1,2)*(-bp);        % from -bp*alpha_dot
B_tau(3) = Minv(1,1);            % from u

% alpha_ddot coefficients
A(4,2) = Minv(2,2)*(-mp*g*hg);
A(4,3) = Minv(2,1)*(-br);
A(4,4) = Minv(2,2)*(-bp);
B_tau(4) = Minv(2,1);

%% ===== Output matrices  =====
C = [ 1 0 0 0 ;
      0 1 0 0 ];

D = zeros(2,1);


sys_tau = ss(A, B_tau, C, D);

disp('A ='); disp(A);
disp('B (tau input) ='); disp(B_tau);
disp('C ='); disp(C);
disp('D ='); disp(D);


B_V = B_tau * (Kt/Rm);                 % input gain to Vm
A_V = A + B_tau * (-(Kt*Km/Rm)) * [0 0 1 0];   % add extra term multiplying theta_dot

sys_V = ss(A_V, B_V, C, D);

disp('--- Using Vm as input ---');
disp('A_V='); disp(A_V);
disp('B_V='); disp(B_V);
\end{lstlisting}
\end{center}


A square-wave voltage input of amplitude $1\,\text{V}$ and frequency $1\,\text{Hz}$ is applied to the system,
\[
V_m(t) = 1 \cdot \mathrm{square}(2\pi \cdot 1\, t),
\]
in order to excite the dynamics while remaining within the validity range of the linear model. The system response is simulated using Simulink model shown below.

\begin{figure}[h!]
\centering
\includegraphics[width=0.85\textwidth]{Photos/sim_down_model.png}
\caption{Simulink simulation model around $\alpha=\pi$ (downward) under a $1\,\text{V}$, $1\,\text{Hz}$ square-wave input.}
\label{fig:sim_down_states}
\end{figure}

\subsection{Simulation Results}

\begin{figure}[h!]
\centering
\includegraphics[width=0.85\textwidth]{Photos/simulationpi.png}
\caption{Simulation around $\alpha=\pi$ (downward): under a square wave input.}
\label{fig:sim_down_states}
\end{figure}


The simulation results exhibit the expected open-loop behavior: the pendulum deviation $\beta(t)$ remains bounded and it performs similar to the theoratical one.



\section{Simulation Around the Upright Equilibrium ($\alpha = 0$)}
Then we simulated around the upright position with zero input, using Simulink with zero initial conditions
\subsection{Simulink Model:}
We similarly built the State Space for the $\alpha$=$\pi$ as in the previuous simulation, but without any input. the output is shown below:
\begin{figure}[h!]
\centering
\includegraphics[width=0.85\textwidth]{Photos/zeroo.jpeg}
\caption{Simulation around $\alpha=0$ under no input.}
\label{fig:sim_up_zero}
\end{figure}


% =========================================================
\section{PID Control: Simulation and Tuning}

In this section, classical feedback control is investigated using two independent controllers: one acting on the pendulum angle $\alpha$ and the other on the rotary arm angle $\theta$. The objective is to improve system stability and transient performance around the upright equilibrium.

% ---------------------------------------------------------
\subsection{Control structure}

Two single-input single-output (SISO) controllers are employed:
\begin{itemize}
  \item A controller regulating the pendulum angle $\alpha$ to maintain balance,
  \item A controller regulating the rotary arm angle $\theta$ to control arm motion.
\end{itemize}

Each controller is implemented as a PID structure. However, through simulation and tuning, it was observed that the integral action led to undesirable effects such as increased overshoot and oscillatory behavior. Consequently, the integral term was removed, and both controllers were implemented as \textbf{PD controllers}.

% ---------------------------------------------------------
\subsection{Simulink implementation}

The control system is implemented in \textsc{Simulink} using the linearized state-space model obtained from \textsc{MATLAB}. The state-space block represents the plant dynamics, while two parallel PD controllers generate the control signals based on the tracking errors in $\alpha$ and $\theta$.

The controller gains were tuned iteratively through simulation to achieve a satisfactory trade-off between stability, response speed, and overshoot. The tuning process was guided by time-domain performance metrics such as settling time and oscillation amplitude.

\begin{figure}[h!]
\centering
\includegraphics[width=0.9\textwidth]{Photos/PIDSimulation.png}
\caption{Simulink model implementing two PD controllers: one for the pendulum angle $\alpha$ and one for the rotary arm angle $\theta$, using the state-space plant defined in MATLAB.}
\label{fig:simulink_pid}
\end{figure}

% ---------------------------------------------------------
\subsection{Simulation results and observations}
\begin{figure}[h!]
\centering
\includegraphics[width=0.9\textwidth]{Photos/2PD.png}
\caption{Simulink model smulation for the 2 PD controllers after tunning.}
\label{fig:simulink_pid_simulation}
\end{figure}
Closed-loop simulations demonstrate a significant improvement in system behavior compared to the open-loop case. The PD controllers are able to stabilize the pendulum near the upright equilibrium and regulate the arm motion and angle $\theta$ with acceptable transient performance. whenever the initial conditions deviate a liitile from 0, the system still performs well.

Also, We note the follwing final constatnts as the starting position when deploying on the hardware:$K_{p\alpha}=50; K_{d\alpha}=20; K_{p\theta}=1;K_{p\theat}=2$. Both controllers used a derivative associated with a LPF with a coefficient $N=10$

The results confirm that derivative action is essential for damping oscillations and improving stability, while integral action is not necessary for this application and may degrade performance. These observations motivate the transition to more systematic control strategies, such as full state feedback, which are explored in the following sections.
% =========================================================

% =========================================================
\section{Deploying, Tuning, and Testing on the Cube Servo}

After validating the PD control strategy through simulation, the controller was deployed on the Cube Servo experimental platform. The objective of this stage was to assess the performance of the controller on real hardware and to account for non-ideal effects not captured in the mathematical model.

% ---------------------------------------------------------
\subsection{Deployment methodology}

The Simulink model developed during the simulation phase was adapted for real-time execution on the Cube Servo system. The state-space plant obtained from \textsc{MATLAB} was replaced by the physical system, while the PD controllers acting on $\alpha$ and $\theta$ were preserved.
\begin{figure}[h!]
\centering
\includegraphics[width=0.9\textwidth]{Photos/deployingPID.png}
\caption{Simulink model used for real-time deployment on the Cube Servo platform, implementing two PD controllers for the pendulum angle $\alpha$ and the rotary arm angle $\theta$.}
\label{fig:cubeservo_simulink}
\end{figure}
Also, we added a comparator that turns on the controller only if the pendulum angle is less than 10\textdegree. Moreover, a saturation block was added to limit the input to the motor to \plusm 5V. Necessarychanges from degrees to radians and vice versa are also done.

The initial controller gains were chosen to be identical to those obtained in simulation. These gains provided a stable starting point; however, due to modeling inaccuracies, sensor noise, friction nonlinearities, and actuator limitations, further tuning was required on the actual hardware.



% ---------------------------------------------------------
\subsection{Experimental tuning}

Controller tuning was performed experimentally using a trial-and-error approach. Starting from the simulation-derived gains, the proportional and derivative coefficients were adjusted incrementally to achieve stable operation and acceptable transient performance.

Through this iterative process, stabilizing gain values were obtained for both control loops:
\begin{itemize}
  \item Pendulum angle controller: $(K_{p,\alpha}, K_{d,\alpha})=(30,7)$
  \item Rotary arm controller: $(K_{p,\theta}, K_{d,\theta})=(2,2)$
\end{itemize}


The final gains achieved reliable stabilization of the pendulum around the upright equilibrium while maintaining smooth arm motion and avoiding excessive oscillations or actuator saturation.

% ---------------------------------------------------------
\subsection{Experimental observations}

The deployed controller successfully stabilized the system under small disturbances and demonstrated behavior consistent with simulation predictions. Minor discrepancies in response speed and oscillation amplitude were observed, primarily due to unmodeled dynamics and measurement noise.

Overall, the experimental results confirm the validity of the proposed PD control strategy and highlight the importance of experimental tuning when transitioning from simulation to real hardware.


% =========================================================
\section{Problems and Limitations of PID Control}

While PID-based control offers simplicity and ease of implementation, it presents several limitations when applied to highly nonlinear and underactuated systems such as the Furuta pendulum. These limitations were observed both in simulation and during experimental deployment on the Cube Servo platform.

First, PID controllers are inherently designed for linear, single-input single-output systems. In the Furuta pendulum, the dynamics are strongly coupled and nonlinear, which makes independent tuning of the $\alpha$ and $\theta$ controllers challenging. Changes in one control loop often affect the behavior of the other, requiring iterative retuning.

Second, the performance of PID control is highly sensitive to operating conditions. Gains tuned around a specific equilibrium may not perform well under larger disturbances or parameter variations. This was particularly evident near the upright equilibrium, where small modeling inaccuracies led to noticeable performance degradation.

Third, derivative action amplifies measurement noise, which can introduce high-frequency oscillations and actuator stress. As a result, excessive derivative gains had to be avoided, limiting achievable performance.

Finally, PID control lacks a systematic design framework for pole placement or robustness guarantees. This makes it difficult to formally predict stability margins and transient behavior, especially when transitioning from simulation to real hardware.

These limitations motivate the use of more advanced control techniques, such as full state feedback, which are addressed in the following sections.
% =========================================================


% =========================================================
\section{Full State Feedback Control}

Full state feedback (FSF) is an advanced control strategy in which the control input is computed using all the states of the system. Unlike classical PID control, which relies primarily on error signals and their derivatives, full state feedback explicitly uses the state-space representation of the system to shape the closed-loop dynamics in a systematic manner.

% ---------------------------------------------------------
\subsection{Concept and motivation}

Consider the linear state-space model
\[
\dot{\bm x} = A\bm x + B u.
\]
In full state feedback control, the input is chosen as
\[
u = -K \bm x,
\]
where $K$ is the state feedback gain vector. The main motivation behind this approach is that, by properly selecting $K$, the closed-loop poles of the system
\[
A_{\text{cl}} = A - BK
\]
can be placed at desired locations in the complex plane, allowing direct control over stability and transient response.

This method provides a more systematic framework for controller design compared to PID control and is widely used in modern control applications.

% ---------------------------------------------------------
\subsection{Pole placement and gain selection}

In this project, several sets of desired closed-loop poles were tested using pole placement techniques in \textsc{MATLAB}. For each pole set, a corresponding feedback gain vector $K$ was computed.

However, it was observed that predicting the system response based solely on pole locations was challenging. Small changes in pole placement often resulted in significantly different time-domain behaviors, making systematic tuning difficult. Additionally, the sensitivity of the system to modeling inaccuracies further complicated the selection of appropriate poles.

% ---------------------------------------------------------
\subsection{Relation to PID control}

In the present case, the full state feedback controller exhibits behavior similar to the previously implemented PD controllers. Since the state vector includes the angles and their time derivatives,
\[
\bm x = [\theta,\ \alpha,\ \dot{\theta},\ \dot{\alpha}]^T,
\]
the control law effectively acts on both the states and their derivatives. As a result, the feedback structure resembles a multivariable PD controller acting directly on the system states.

This similarity explains why the performance improvement over the tuned PD controllers was limited in this specific application.
\begin{figure}[h!]
\centering
\includegraphics[width=0.9\textwidth]{Photos/fullstate.png}
\caption{Simulink model implementing full state feedback control using the representation of the Furuta pendulum.}
\label{fig:simulink_fullstate}
\end{figure}
\subsection{State measurement and estimation}

While the angular positions $\theta$ and $\alpha$ are directly measured, the angular velocities $\dot{\theta}$ and $\dot{\alpha}$ are not available as direct measurements. Therefore, the velocities were estimated by numerically differentiating the measured angles and applying a low-pass filter (LPF) to reduce noise amplification.

Although this approach provides reasonable estimates, it introduces phase lag and measurement noise, which negatively affect controller performance and stability margins.

% ---------------------------------------------------------
\subsection{Simulink implementation and deployment}

The full state feedback controller was implemented in \textsc{Simulink} using the linearized state-space model. The computed gain vector $K$ was applied to the state vector to generate the control input.



Deployment on the Cube Servo platform did not result in stable operation. The controller exhibited sensitivity to noise, unmodeled dynamics, and state estimation errors, which ultimately prevented reliable stabilization of the system.

These results highlight the practical challenges associated with full state feedback implementation and emphasize the importance of accurate state measurement and robust estimation techniques, such as observers or Kalman filters, LQR, and Fuzzy Logic, which were beyond the scope of this project.
% =========================================================


\section{Conclusion}

This project presented a comprehensive modeling, analysis, and control study of the Furuta pendulum system. Starting from first principles, the nonlinear dynamics were derived using Lagrangian mechanics, resulting in a coupled set of equations that accurately capture the interaction between the rotary arm and the pendulum link. These equations formed the basis for all subsequent analysis and controller design.

Linearization was performed around both the downward equilibrium ($\alpha=\pi$) and the upright equilibrium ($\alpha=0$), revealing fundamentally different system behaviors. The downward configuration was shown to be naturally stable, while the upright configuration was confirmed to be inherently unstable and requiring active control. State-space representations were derived explicitly and used to analyze system dynamics, simulate responses, and implement control strategies.

Several control approaches were explored. Open-loop simulations validated the analytical models and highlighted the instability of the upright equilibrium. Classical PID control was then investigated, and through simulation and experimental tuning, PD controllers were found to provide reliable stabilization of the system around the upright position. The successful deployment of the PD controllers on the Cube Servo platform demonstrated the practicality of classical control techniques when carefully tuned and experimentally validated.

Full state feedback control was also studied as a more advanced control strategy. While pole placement techniques were successfully applied in computation via MATLAB, deployment on real hardware proved challenging due to sensitivity to noise, unmodeled dynamics, and state estimation errors. In practice, the full state feedback implementation exhibited behavior similar to the PD controllers, as both relied on feedback from the states and their derivatives. These results emphasize the importance of accurate state measurement and robust estimation methods when implementing advanced control techniques.

Overall, this project effectively bridged theoretical control concepts with real-world implementation. It highlighted the challenges associated with nonlinear systems, equilibrium-dependent behavior, and hardware constraints, while reinforcing key ideas in modeling, linearization, state-space analysis, and feedback control. The Furuta pendulum experiment proved to be a valuable learning platform for understanding both the power and limitations of control system design in practice.

\end{document}
